\documentclass[a4paper,11pt]{book} 

%=== PACKAGES ==============================================
\usepackage{graphicx}
\usepackage[english]{babel}
\usepackage{color}      % for colored text
\usepackage{xspace}     % for explicit spaces in abbreviations (\xspace)
\usepackage{moreverb}   % more verbatim: \begin{boxedverbatim} for source code
\usepackage{verbatim}   % for plain text
\usepackage{algorithm}  % for algorithms
\usepackage{algorithmic}% for algorithms
\usepackage{multirow}   % for multi-rows in tables
\usepackage{cite}       % will become: [2], [5]--[7]  (using cite.sty)
\usepackage{float}
\usepackage{listings}
\usepackage{array}
\usepackage{makeidx,shortvrb,latexsym}
\usepackage[latin1]{inputenc}
\usepackage{a4wide}
\usepackage{amsfonts}
\usepackage{amsmath}
\usepackage{amssymb}
%\usepackage{MnSymbol}
\usepackage{fancyhdr}
\usepackage{subfigure}
\usepackage{psfrag}
\usepackage{epsfig}
\usepackage{caption}
\usepackage{tocvsec2}
%\usepackage[draft]{pdfpages}
\usepackage{colortbl}
\usepackage[T1]{fontenc}

%% =============================================
%% LaTeX Macros 
%% =============================================
%%
%%
%% =============================================

%% -- References 
\renewcommand{\eqref}[1]{(\ref{#1})}
%\renewcommand{\eqref}[1]{Eq.~\ref{#1}}
\newcommand{\secref}[1]{Section~\ref{#1}}
\newcommand{\chapref}[1]{Chapter~\ref{#1}}
\newcommand{\figref}[1]{Fig.~\ref{#1}}
\newcommand{\tblref}[1]{Table~\ref{#1}}

%% -- soft information specific
\newcommand{\lambdaML}[1]{\ensuremath{\lambda^\mathrm{ML}_{#1}}}
\newcommand{\lambdaMLB}[1]{\ensuremath{\lambda^{\overline{\mathrm{ML}}}_{#1}}}
\newcommand{\xML}[1]{\ensuremath{x^\mathrm{ML}_{#1}}}
\newcommand{\xMLB}[1]{\ensuremath{\overline{x^\mathrm{ML}_{#1}}}}

%% -- Probabilities
\renewcommand{\Pr}[1]{\ensuremath{\mathrm{Pr}\left[#1\right]}}
\newcommand{\PDF}[1]{\ensuremath{\mathrm{p}\left(#1\right)}}
\newcommand{\Ex}{\ensuremath{\mathbb{E}}}
\newcommand{\var}[1]{\ensuremath{\mathrm{var}\left[#1\right]}}

%% -- Operators
\DeclareMathOperator*{\argmin}{arg\;min}
\DeclareMathOperator*{\argmax}{arg\;max}
\DeclareMathOperator*{\arc}{arc}
\DeclareMathOperator*{\Tr}{Tr}

%% -- Other Math
\newcommand{\dhat}{\hat{d}}
\newcommand{\xihat}{\hat{\xi}}
\newcommand{\indep}{{\;\bot\!\!\!\!\!\!\bot\;\,}}

%% -- Various commands
\newcommand{\slice}[1]{\ensuremath{\mathrm{Q}\left(#1\right)}}
\newcommand{\COrder}[1]{\ensuremath{\mathcal{O}(#1)}}
\newcommand{\vecnorm}[1]{\ensuremath{\lVert#1\rVert}}
\newcommand{\adj}[1]{\ensuremath{\mathrm{adj}\,#1}}
\newcommand{\card}[1]{\ensuremath{|#1|}}
\newcommand{\abs}[1]{\ensuremath{\left|#1\right|}}
\newcommand{\expect}[1]{\ensuremath{\mathcal{E}\{#1\}}}
\newcommand{\NP}{\ensuremath{\mathcal{NP}}}
\newcommand{\eigv}[1]{\ensuremath{\lambda(#1)}}
\newcommand{\eigvmin}[1]{\ensuremath{\lambda_{\mathrm{min}}(#1)}}
\newcommand{\eigvmax}[1]{\ensuremath{\lambda_{\mathrm{max}}(#1)}}
\newcommand{\Cdim}[2]{\ensuremath{\in\mathbb{C}^{#1\times#2}}}
\newcommand{\Rdim}[2]{\ensuremath{\in\mathbb{R}^{#1\times#2}}}
\newcommand{\with}{\ensuremath{\quad\mathrm{with}\quad}}
\newcommand{\diag}[1]{\ensuremath{\mathrm{diag}\left\lbrace#1\right\rbrace}}
\newcommand{\atan}{\ensuremath{\mathrm{atan}}}
\newcommand{\Real}{\ensuremath{\mathbb{R}}}
\newcommand{\Complex}{\ensuremath{\mathbb{C}}}
\newcommand{\cmplx}{\ensuremath{\mathsf{i}}}
\newcommand{\constant}{\ensuremath{\mathrm{const.}}}
\newcommand{\floor}[1]{\ensuremath{\lfloor#1\rfloor}}
\newcommand{\ceil}[1]{\ensuremath{\lceil#1\rceil}}
\newcommand{\round}[1]{\ensuremath{\mathrm{int}\left(#1\right)}}
\newcommand{\ltwo}{\ensuremath{\ell^{2}}}
\newcommand{\lone}{\ensuremath{\ell^{1}}}
\newcommand{\linfty}{\ensuremath{\ell^{\infty}}}

%% -- sans serif
\newcommand{\sfA}{\ensuremath{\mathsf{A}}}
\newcommand{\sfB}{\ensuremath{\mathsf{B}}}
\newcommand{\sfC}{\ensuremath{\mathsf{C}}}
\newcommand{\sfD}{\ensuremath{\mathsf{D}}}
\newcommand{\sfE}{\ensuremath{\mathsf{E}}}
\newcommand{\sfF}{\ensuremath{\mathsf{F}}}
\newcommand{\sfG}{\ensuremath{\mathsf{G}}}
\newcommand{\sfH}{\ensuremath{\mathsf{H}}}
\newcommand{\sfI}{\ensuremath{\mathsf{I}}}
\newcommand{\sfJ}{\ensuremath{\mathsf{J}}}
\newcommand{\sfK}{\ensuremath{\mathsf{K}}}
\newcommand{\sfL}{\ensuremath{\mathsf{L}}}
\newcommand{\sfM}{\ensuremath{\mathsf{M}}}
\newcommand{\sfN}{\ensuremath{\mathsf{N}}}
\newcommand{\sfO}{\ensuremath{\mathsf{O}}}
\newcommand{\sfP}{\ensuremath{\mathsf{P}}}
\newcommand{\sfQ}{\ensuremath{\mathsf{Q}}}
\newcommand{\sfR}{\ensuremath{\mathsf{R}}}
\newcommand{\sfS}{\ensuremath{\mathsf{S}}}
\newcommand{\sfT}{\ensuremath{\mathsf{T}}}
\newcommand{\sfU}{\ensuremath{\mathsf{U}}}
\newcommand{\sfV}{\ensuremath{\mathsf{V}}}
\newcommand{\sfW}{\ensuremath{\mathsf{W}}}
\newcommand{\sfX}{\ensuremath{\mathsf{X}}}
\newcommand{\sfY}{\ensuremath{\mathsf{Y}}}
\newcommand{\sfZ}{\ensuremath{\mathsf{Z}}}

%% -- Sets
\newcommand{\setA}{\ensuremath{\mathcal{A}}}
\newcommand{\setB}{\ensuremath{\mathcal{B}}}
\newcommand{\setC}{\ensuremath{\mathcal{C}}}
\newcommand{\setD}{\ensuremath{\mathcal{D}}}
\newcommand{\setE}{\ensuremath{\mathcal{E}}}
\newcommand{\setF}{\ensuremath{\mathcal{F}}}
\newcommand{\setG}{\ensuremath{\mathcal{G}}}
\newcommand{\setH}{\ensuremath{\mathcal{H}}}
\newcommand{\setI}{\ensuremath{\mathcal{I}}}
\newcommand{\setJ}{\ensuremath{\mathcal{J}}}
\newcommand{\setK}{\ensuremath{\mathcal{K}}}
\newcommand{\setL}{\ensuremath{\mathcal{L}}}
\newcommand{\setM}{\ensuremath{\mathcal{M}}}
\newcommand{\setN}{\ensuremath{\mathcal{N}}}
\newcommand{\setO}{\ensuremath{\mathcal{O}}}
\newcommand{\setP}{\ensuremath{\mathcal{P}}}
\newcommand{\setQ}{\ensuremath{\mathcal{Q}}}
\newcommand{\setR}{\ensuremath{\mathcal{R}}}
\newcommand{\setS}{\ensuremath{\mathcal{S}}}
\newcommand{\setT}{\ensuremath{\mathcal{T}}}
\newcommand{\setU}{\ensuremath{\mathcal{U}}}
\newcommand{\setV}{\ensuremath{\mathcal{V}}}
\newcommand{\setW}{\ensuremath{\mathcal{W}}}
\newcommand{\setX}{\ensuremath{\mathcal{X}}}
\newcommand{\setY}{\ensuremath{\mathcal{Y}}}
\newcommand{\setZ}{\ensuremath{\mathcal{Z}}}

%% -- Lower case letters (vectors)
\newcommand{\bma}{\ensuremath{\mathbf{a}}}
\newcommand{\bmb}{\ensuremath{\mathbf{b}}}
\newcommand{\bmc}{\ensuremath{\mathbf{c}}}
\newcommand{\bmd}{\ensuremath{\mathbf{d}}}
\newcommand{\bme}{\ensuremath{\mathbf{e}}}
\newcommand{\bmf}{\ensuremath{\mathbf{f}}}
\newcommand{\bmg}{\ensuremath{\mathbf{g}}}
\newcommand{\bmh}{\ensuremath{\mathbf{h}}}
\newcommand{\bmi}{\ensuremath{\mathbf{i}}}
\newcommand{\bmj}{\ensuremath{\mathbf{j}}}
\newcommand{\bmk}{\ensuremath{\mathbf{k}}}
\newcommand{\bml}{\ensuremath{\mathbf{l}}}
\newcommand{\bmm}{\ensuremath{\mathbf{m}}}
\newcommand{\bmn}{\ensuremath{\mathbf{n}}}
\newcommand{\bmo}{\ensuremath{\mathbf{o}}}
\newcommand{\bmp}{\ensuremath{\mathbf{p}}}
\newcommand{\bmq}{\ensuremath{\mathbf{q}}}
\newcommand{\bmr}{\ensuremath{\mathbf{r}}}
\newcommand{\bms}{\ensuremath{\mathbf{s}}}
\newcommand{\bmt}{\ensuremath{\mathbf{t}}}
\newcommand{\bmu}{\ensuremath{\mathbf{u}}}
\newcommand{\bmv}{\ensuremath{\mathbf{v}}}
\newcommand{\bmw}{\ensuremath{\mathbf{w}}}
\newcommand{\bmx}{\ensuremath{\mathbf{x}}}
\newcommand{\bmy}{\ensuremath{\mathbf{y}}}
\newcommand{\bmz}{\ensuremath{\mathbf{z}}}

%% -- Upper case letters (Matrices)
\newcommand{\bA}{\ensuremath{\mathbf{A}}}
\newcommand{\bB}{\ensuremath{\mathbf{B}}}
\newcommand{\bC}{\ensuremath{\mathbf{C}}}
\newcommand{\bD}{\ensuremath{\mathbf{D}}}
\newcommand{\bE}{\ensuremath{\mathbf{E}}}
\newcommand{\bF}{\ensuremath{\mathbf{F}}}
\newcommand{\bG}{\ensuremath{\mathbf{G}}}
\newcommand{\bH}{\ensuremath{\mathbf{H}}}
\newcommand{\bI}{\ensuremath{\mathbf{I}}}
\newcommand{\bJ}{\ensuremath{\mathbf{J}}}
\newcommand{\bK}{\ensuremath{\mathbf{K}}}
\newcommand{\bL}{\ensuremath{\mathbf{L}}}
\newcommand{\bM}{\ensuremath{\mathbf{M}}}
\newcommand{\bN}{\ensuremath{\mathbf{N}}}
\newcommand{\bO}{\ensuremath{\mathbf{O}}}
\newcommand{\bP}{\ensuremath{\mathbf{P}}}
\newcommand{\bQ}{\ensuremath{\mathbf{Q}}}
\newcommand{\bR}{\ensuremath{\mathbf{R}}}
\newcommand{\bS}{\ensuremath{\mathbf{S}}}
\newcommand{\bT}{\ensuremath{\mathbf{T}}}
\newcommand{\bU}{\ensuremath{\mathbf{U}}}
\newcommand{\bV}{\ensuremath{\mathbf{V}}}
\newcommand{\bW}{\ensuremath{\mathbf{W}}}
\newcommand{\bX}{\ensuremath{\mathbf{X}}}
\newcommand{\bY}{\ensuremath{\mathbf{Y}}}
\newcommand{\bZ}{\ensuremath{\mathbf{Z}}}
\newcommand{\bZero}{\ensuremath{\mathbf{0}}}
\newcommand{\bOne}{\ensuremath{\mathbf{1}}}
\newcommand{\bPsi}{\ensuremath{\mathbf{\Psi}}}
\newcommand{\bTheta}{\ensuremath{\mathbf{\Theta}}}
\newcommand{\bDelta}{\ensuremath{\mathbf{\Delta}}}
\newcommand{\bLambda}{\ensuremath{\mathbf{\Lambda}}}

%% -- Vectors with hat
\newcommand{\bahat}{\ensuremath{\hat{\bma}}}
\newcommand{\bbhat}{\ensuremath{\hat{\bmb}}}
\newcommand{\bchat}{\ensuremath{\hat{\bmc}}}
\newcommand{\bdhat}{\ensuremath{\hat{\bmd}}}
\newcommand{\behat}{\ensuremath{\hat{\bme}}}
\newcommand{\bfhat}{\ensuremath{\hat{\bmf}}}
\newcommand{\bghat}{\ensuremath{\hat{\bmg}}}
\newcommand{\bhhat}{\ensuremath{\hat{\bmh}}}
\newcommand{\bihat}{\ensuremath{\hat{\bmi}}}
\newcommand{\bjhat}{\ensuremath{\hat{\bmj}}}
\newcommand{\bkhat}{\ensuremath{\hat{\bmk}}}
\newcommand{\blhat}{\ensuremath{\hat{\bml}}}
\newcommand{\bmhat}{\ensuremath{\hat{\bmm}}}
\newcommand{\bnhat}{\ensuremath{\hat{\bmn}}}
\newcommand{\bohat}{\ensuremath{\hat{\bmo}}}
\newcommand{\bphat}{\ensuremath{\hat{\bmp}}}
\newcommand{\bqhat}{\ensuremath{\hat{\bmq}}}
\newcommand{\brhat}{\ensuremath{\hat{\bmr}}}
\newcommand{\bshat}{\ensuremath{\hat{\bms}}}
\newcommand{\bthat}{\ensuremath{\hat{\bmt}}}
\newcommand{\buhat}{\ensuremath{\hat{\bmu}}}
\newcommand{\bvhat}{\ensuremath{\hat{\bmv}}}
\newcommand{\bwhat}{\ensuremath{\hat{\bmw}}}
\newcommand{\bxhat}{\ensuremath{\hat{\bmx}}}
\newcommand{\byhat}{\ensuremath{\hat{\bmy}}}
\newcommand{\bzhat}{\ensuremath{\hat{\bmz}}}

%% -- Vectors with tilde
\newcommand{\batilde}{\ensuremath{\tilde{\bma}}}
\newcommand{\bbtilde}{\ensuremath{\tilde{\bmb}}}
\newcommand{\bctilde}{\ensuremath{\tilde{\bmc}}}
\newcommand{\bdtilde}{\ensuremath{\tilde{\bmd}}}
\newcommand{\betilde}{\ensuremath{\tilde{\bme}}}
\newcommand{\bftilde}{\ensuremath{\tilde{\bmf}}}
\newcommand{\bgtilde}{\ensuremath{\tilde{\bmg}}}
\newcommand{\bhtilde}{\ensuremath{\tilde{\bmh}}}
\newcommand{\bitilde}{\ensuremath{\tilde{\bmi}}}
\newcommand{\bjtilde}{\ensuremath{\tilde{\bmj}}}
\newcommand{\bktilde}{\ensuremath{\tilde{\bmk}}}
\newcommand{\bltilde}{\ensuremath{\tilde{\bml}}}
\newcommand{\bmtilde}{\ensuremath{\tilde{\bmm}}}
\newcommand{\bntilde}{\ensuremath{\tilde{\bmn}}}
\newcommand{\botilde}{\ensuremath{\tilde{\bmo}}}
\newcommand{\bptilde}{\ensuremath{\tilde{\bmp}}}
\newcommand{\bqtilde}{\ensuremath{\tilde{\bmq}}}
\newcommand{\brtilde}{\ensuremath{\tilde{\bmr}}}
\newcommand{\bstilde}{\ensuremath{\tilde{\bms}}}
\newcommand{\bttilde}{\ensuremath{\tilde{\bmt}}}
\newcommand{\butilde}{\ensuremath{\tilde{\bmu}}}
\newcommand{\bvtilde}{\ensuremath{\tilde{\bmv}}}
\newcommand{\bwtilde}{\ensuremath{\tilde{\bmw}}}
\newcommand{\bxtilde}{\ensuremath{\tilde{\bmx}}}
\newcommand{\bytilde}{\ensuremath{\tilde{\bmy}}}
\newcommand{\bztilde}{\ensuremath{\tilde{\bmz}}}

%% -- Lattices
\newcommand{\lLAMBDA}{\ensuremath{\mathbf{\Lambda}}}

%% -- MIMO specific macros
\newcommand{\MT}{\ensuremath{M_{T}}}
\newcommand{\MR}{\ensuremath{M_{R}}}
\newcommand{\SNR}{\ensuremath{\mathsf{SNR}}}
\newcommand{\SQNR}{\ensuremath{\mathrm{SQNR}}}
\newcommand{\SINR}{\ensuremath{\textsf{SINR}}}
\newcommand{\const}{\ensuremath{\setO}}
\newcommand{\mimoconst}{\ensuremath{\setO^{\MT}}}
\newcommand{\modorder}{\ensuremath{Q}}
\newcommand{\bHH}{\ensuremath{\bH^{H}}}
\newcommand{\bHbar}{\ensuremath{\bar{\bH}}}
\newcommand{\bQbar}{\ensuremath{\bar{\bQ}}}
\newcommand{\bRbar}{\ensuremath{\bar{\bR}}}
\newcommand{\bQa}{\ensuremath{\bQ_{\alpha}}}
\newcommand{\bRa}{\ensuremath{\bR_{\alpha}}}
\newcommand{\bIbar}{\ensuremath{\bar{\bI}}}
\newcommand{\bmybar}{\ensuremath{\bar{\bmy}}}
\newcommand{\bHi}{\ensuremath{\bH^{\dag}}}
\newcommand{\Pn}{\ensuremath{\sigma^{2}}}
\newcommand{\No}{\ensuremath{N_{o}}}
\newcommand{\VBLAST}{\ensuremath{\mbox{V-BLAST}}}
\newcommand{\Rate}{\ensuremath{R}}
\newcommand{\bsML}{\ensuremath{\bms^{\mathrm{ML}}}}
\newcommand{\bsZF}{\ensuremath{\bms^{\mathrm{ZF}}}}
\newcommand{\bsMMSE}{\ensuremath{\bms^{\mathrm{MMSE}}}}

%% -- VLSI things
\newcommand{\tclk}{\ensuremath{t_\mathrm{clk}}}
\newcommand{\fclk}{\ensuremath{f_\mathrm{clk}}}

%% -- Complexities
\newcommand{\CGr}{\ensuremath{C_{\mathrm{GR}}}}
\newcommand{\CVec}{\ensuremath{C_{\mathrm{Vec}}}}
\newcommand{\CRot}{\ensuremath{C_{\mathrm{Rot}}}}
\newcommand{\CMult}{\ensuremath{C_{\mathrm{Mult}}}}
\newcommand{\COMult}{\ensuremath{C_{\mathrm{OptMult}}}}
\newcommand{\CDiv}{\ensuremath{C_{\mathrm{Div}}}}
\newcommand{\CSqrt}{\ensuremath{C_{\sqrt{\cdot}}}}

%% -- NON-Math macros
\newcommand{\UMC}{$0.25\mu\mathrm{m}$ 1P/5M technology}

%% -- For NOTATION
\def\addnotation #1: #2{\parbox[t]{1.25cm}{$#1$\dotfill} \parbox[t]{3.50in}{#2} \vspace{0.0625cm} \\}
\def\addsymbol #1: #2{\parbox[t]{2cm}{$#1$} \parbox[t]{3in}{#2} \vspace{0.0625cm} \\}
\def\addacronym #1: #2{\parbox[t]{2cm}{\bf{#1}\dotfill} \parbox[t]{2.8in}{#2} \vspace{0.0625cm} \\}


\definecolor{lgray}{gray}{0.9}
\definecolor{dblue}{rgb}{0,0.08,0.7}
%\definecolor{dblue}{rgb}{0.1,0.1,0.45}
\definecolor{lred}{rgb}{1.0,0.8,0.8} 
\definecolor{lgreen}{rgb}{0.8,1.0,0.8} 

%=== Colored Hyperlinks ======================================
\usepackage[ % IF PDF
  colorlinks = true,        % links are colored
  citecolor  = dblue,        % color of cite links
  pagecolor  = dblue,        % color of page links
  linkcolor  = dblue,        % color of hyperref links
  menucolor  = dblue,        % color of Acrobat Reader menu buttons
  urlcolor   = dblue,        % color of urls \url{}
  bookmarks  = true, hyperindex=true, breaklinks=true,
  baseurl={http://www.nari.ee.ethz.ch},
  pdftitle={Distributed Multiuser Multihop Networks},             
  pdfauthor={Raphael Rolny},             
  pdfsubject={xxx},
  pdfkeywords={xxx}
]{hyperref}
    % IF LaTeX
%\usepackage[colorlinks=false]{hyperref}

%=== Theorems ================================================
\usepackage{amsthm}
\newtheorem{thm}{Theorem}[chapter]
\newtheorem{exmpl}{Example}[chapter]
\newtheorem{propos}{Proposition}[chapter]
\newtheorem{defi}{Definition}[chapter]
\newtheorem{corol}{Corollary}[chapter]
\newtheorem{lemma}{Lemma}[chapter]

%== for algorithms ============================================
\floatstyle{ruled}
\newfloat{algorithm}{tb}{loa}  % loa: auxiliary file for list of algorithms
\floatname{algorithm}{Algorithm}
\newcommand{\theHalgorithm}{\arabic{algorithm}}

%=== Index ====================================================
%\usepackage{makeidx} 
%\makeindex 

%=== Fancy header =============================================
\pagestyle{fancyplain}
% with this we ensure that the chapter and section headings are in lowercase
\fancyhf{} % delete current setting for header and footer
\fancyhead[LE,RO]{\thepage}
\fancyhead[LO]{\rightmark}
\fancyhead[RE]{\leftmark}
\addtolength{\headheight}{1.6pt} % make space for the rule
\renewcommand{\headrulewidth}{0.5pt}
\fancypagestyle{plain}{
   \fancyhead{}  % get rid of headers on plain pages
   \renewcommand{\headrulewidth}{0pt} % and the line
}


% --- Definition of short-cuts/macros ------------------------------
%%% =============================================
%% LaTeX Macros 
%% =============================================
%%
%%
%% =============================================

%% -- References 
\renewcommand{\eqref}[1]{(\ref{#1})}
%\renewcommand{\eqref}[1]{Eq.~\ref{#1}}
\newcommand{\secref}[1]{Section~\ref{#1}}
\newcommand{\chapref}[1]{Chapter~\ref{#1}}
\newcommand{\figref}[1]{Fig.~\ref{#1}}
\newcommand{\tblref}[1]{Table~\ref{#1}}

%% -- soft information specific
\newcommand{\lambdaML}[1]{\ensuremath{\lambda^\mathrm{ML}_{#1}}}
\newcommand{\lambdaMLB}[1]{\ensuremath{\lambda^{\overline{\mathrm{ML}}}_{#1}}}
\newcommand{\xML}[1]{\ensuremath{x^\mathrm{ML}_{#1}}}
\newcommand{\xMLB}[1]{\ensuremath{\overline{x^\mathrm{ML}_{#1}}}}

%% -- Probabilities
\renewcommand{\Pr}[1]{\ensuremath{\mathrm{Pr}\left[#1\right]}}
\newcommand{\PDF}[1]{\ensuremath{\mathrm{p}\left(#1\right)}}
\newcommand{\Ex}{\ensuremath{\mathbb{E}}}
\newcommand{\var}[1]{\ensuremath{\mathrm{var}\left[#1\right]}}

%% -- Operators
\DeclareMathOperator*{\argmin}{arg\;min}
\DeclareMathOperator*{\argmax}{arg\;max}
\DeclareMathOperator*{\arc}{arc}
\DeclareMathOperator*{\Tr}{Tr}

%% -- Other Math
\newcommand{\dhat}{\hat{d}}
\newcommand{\xihat}{\hat{\xi}}
\newcommand{\indep}{{\;\bot\!\!\!\!\!\!\bot\;\,}}

%% -- Various commands
\newcommand{\slice}[1]{\ensuremath{\mathrm{Q}\left(#1\right)}}
\newcommand{\COrder}[1]{\ensuremath{\mathcal{O}(#1)}}
\newcommand{\vecnorm}[1]{\ensuremath{\lVert#1\rVert}}
\newcommand{\adj}[1]{\ensuremath{\mathrm{adj}\,#1}}
\newcommand{\card}[1]{\ensuremath{|#1|}}
\newcommand{\abs}[1]{\ensuremath{\left|#1\right|}}
\newcommand{\expect}[1]{\ensuremath{\mathcal{E}\{#1\}}}
\newcommand{\NP}{\ensuremath{\mathcal{NP}}}
\newcommand{\eigv}[1]{\ensuremath{\lambda(#1)}}
\newcommand{\eigvmin}[1]{\ensuremath{\lambda_{\mathrm{min}}(#1)}}
\newcommand{\eigvmax}[1]{\ensuremath{\lambda_{\mathrm{max}}(#1)}}
\newcommand{\Cdim}[2]{\ensuremath{\in\mathbb{C}^{#1\times#2}}}
\newcommand{\Rdim}[2]{\ensuremath{\in\mathbb{R}^{#1\times#2}}}
\newcommand{\with}{\ensuremath{\quad\mathrm{with}\quad}}
\newcommand{\diag}[1]{\ensuremath{\mathrm{diag}\left\lbrace#1\right\rbrace}}
\newcommand{\atan}{\ensuremath{\mathrm{atan}}}
\newcommand{\Real}{\ensuremath{\mathbb{R}}}
\newcommand{\Complex}{\ensuremath{\mathbb{C}}}
\newcommand{\cmplx}{\ensuremath{\mathsf{i}}}
\newcommand{\constant}{\ensuremath{\mathrm{const.}}}
\newcommand{\floor}[1]{\ensuremath{\lfloor#1\rfloor}}
\newcommand{\ceil}[1]{\ensuremath{\lceil#1\rceil}}
\newcommand{\round}[1]{\ensuremath{\mathrm{int}\left(#1\right)}}
\newcommand{\ltwo}{\ensuremath{\ell^{2}}}
\newcommand{\lone}{\ensuremath{\ell^{1}}}
\newcommand{\linfty}{\ensuremath{\ell^{\infty}}}

%% -- sans serif
\newcommand{\sfA}{\ensuremath{\mathsf{A}}}
\newcommand{\sfB}{\ensuremath{\mathsf{B}}}
\newcommand{\sfC}{\ensuremath{\mathsf{C}}}
\newcommand{\sfD}{\ensuremath{\mathsf{D}}}
\newcommand{\sfE}{\ensuremath{\mathsf{E}}}
\newcommand{\sfF}{\ensuremath{\mathsf{F}}}
\newcommand{\sfG}{\ensuremath{\mathsf{G}}}
\newcommand{\sfH}{\ensuremath{\mathsf{H}}}
\newcommand{\sfI}{\ensuremath{\mathsf{I}}}
\newcommand{\sfJ}{\ensuremath{\mathsf{J}}}
\newcommand{\sfK}{\ensuremath{\mathsf{K}}}
\newcommand{\sfL}{\ensuremath{\mathsf{L}}}
\newcommand{\sfM}{\ensuremath{\mathsf{M}}}
\newcommand{\sfN}{\ensuremath{\mathsf{N}}}
\newcommand{\sfO}{\ensuremath{\mathsf{O}}}
\newcommand{\sfP}{\ensuremath{\mathsf{P}}}
\newcommand{\sfQ}{\ensuremath{\mathsf{Q}}}
\newcommand{\sfR}{\ensuremath{\mathsf{R}}}
\newcommand{\sfS}{\ensuremath{\mathsf{S}}}
\newcommand{\sfT}{\ensuremath{\mathsf{T}}}
\newcommand{\sfU}{\ensuremath{\mathsf{U}}}
\newcommand{\sfV}{\ensuremath{\mathsf{V}}}
\newcommand{\sfW}{\ensuremath{\mathsf{W}}}
\newcommand{\sfX}{\ensuremath{\mathsf{X}}}
\newcommand{\sfY}{\ensuremath{\mathsf{Y}}}
\newcommand{\sfZ}{\ensuremath{\mathsf{Z}}}

%% -- Sets
\newcommand{\setA}{\ensuremath{\mathcal{A}}}
\newcommand{\setB}{\ensuremath{\mathcal{B}}}
\newcommand{\setC}{\ensuremath{\mathcal{C}}}
\newcommand{\setD}{\ensuremath{\mathcal{D}}}
\newcommand{\setE}{\ensuremath{\mathcal{E}}}
\newcommand{\setF}{\ensuremath{\mathcal{F}}}
\newcommand{\setG}{\ensuremath{\mathcal{G}}}
\newcommand{\setH}{\ensuremath{\mathcal{H}}}
\newcommand{\setI}{\ensuremath{\mathcal{I}}}
\newcommand{\setJ}{\ensuremath{\mathcal{J}}}
\newcommand{\setK}{\ensuremath{\mathcal{K}}}
\newcommand{\setL}{\ensuremath{\mathcal{L}}}
\newcommand{\setM}{\ensuremath{\mathcal{M}}}
\newcommand{\setN}{\ensuremath{\mathcal{N}}}
\newcommand{\setO}{\ensuremath{\mathcal{O}}}
\newcommand{\setP}{\ensuremath{\mathcal{P}}}
\newcommand{\setQ}{\ensuremath{\mathcal{Q}}}
\newcommand{\setR}{\ensuremath{\mathcal{R}}}
\newcommand{\setS}{\ensuremath{\mathcal{S}}}
\newcommand{\setT}{\ensuremath{\mathcal{T}}}
\newcommand{\setU}{\ensuremath{\mathcal{U}}}
\newcommand{\setV}{\ensuremath{\mathcal{V}}}
\newcommand{\setW}{\ensuremath{\mathcal{W}}}
\newcommand{\setX}{\ensuremath{\mathcal{X}}}
\newcommand{\setY}{\ensuremath{\mathcal{Y}}}
\newcommand{\setZ}{\ensuremath{\mathcal{Z}}}

%% -- Lower case letters (vectors)
\newcommand{\bma}{\ensuremath{\mathbf{a}}}
\newcommand{\bmb}{\ensuremath{\mathbf{b}}}
\newcommand{\bmc}{\ensuremath{\mathbf{c}}}
\newcommand{\bmd}{\ensuremath{\mathbf{d}}}
\newcommand{\bme}{\ensuremath{\mathbf{e}}}
\newcommand{\bmf}{\ensuremath{\mathbf{f}}}
\newcommand{\bmg}{\ensuremath{\mathbf{g}}}
\newcommand{\bmh}{\ensuremath{\mathbf{h}}}
\newcommand{\bmi}{\ensuremath{\mathbf{i}}}
\newcommand{\bmj}{\ensuremath{\mathbf{j}}}
\newcommand{\bmk}{\ensuremath{\mathbf{k}}}
\newcommand{\bml}{\ensuremath{\mathbf{l}}}
\newcommand{\bmm}{\ensuremath{\mathbf{m}}}
\newcommand{\bmn}{\ensuremath{\mathbf{n}}}
\newcommand{\bmo}{\ensuremath{\mathbf{o}}}
\newcommand{\bmp}{\ensuremath{\mathbf{p}}}
\newcommand{\bmq}{\ensuremath{\mathbf{q}}}
\newcommand{\bmr}{\ensuremath{\mathbf{r}}}
\newcommand{\bms}{\ensuremath{\mathbf{s}}}
\newcommand{\bmt}{\ensuremath{\mathbf{t}}}
\newcommand{\bmu}{\ensuremath{\mathbf{u}}}
\newcommand{\bmv}{\ensuremath{\mathbf{v}}}
\newcommand{\bmw}{\ensuremath{\mathbf{w}}}
\newcommand{\bmx}{\ensuremath{\mathbf{x}}}
\newcommand{\bmy}{\ensuremath{\mathbf{y}}}
\newcommand{\bmz}{\ensuremath{\mathbf{z}}}

%% -- Upper case letters (Matrices)
\newcommand{\bA}{\ensuremath{\mathbf{A}}}
\newcommand{\bB}{\ensuremath{\mathbf{B}}}
\newcommand{\bC}{\ensuremath{\mathbf{C}}}
\newcommand{\bD}{\ensuremath{\mathbf{D}}}
\newcommand{\bE}{\ensuremath{\mathbf{E}}}
\newcommand{\bF}{\ensuremath{\mathbf{F}}}
\newcommand{\bG}{\ensuremath{\mathbf{G}}}
\newcommand{\bH}{\ensuremath{\mathbf{H}}}
\newcommand{\bI}{\ensuremath{\mathbf{I}}}
\newcommand{\bJ}{\ensuremath{\mathbf{J}}}
\newcommand{\bK}{\ensuremath{\mathbf{K}}}
\newcommand{\bL}{\ensuremath{\mathbf{L}}}
\newcommand{\bM}{\ensuremath{\mathbf{M}}}
\newcommand{\bN}{\ensuremath{\mathbf{N}}}
\newcommand{\bO}{\ensuremath{\mathbf{O}}}
\newcommand{\bP}{\ensuremath{\mathbf{P}}}
\newcommand{\bQ}{\ensuremath{\mathbf{Q}}}
\newcommand{\bR}{\ensuremath{\mathbf{R}}}
\newcommand{\bS}{\ensuremath{\mathbf{S}}}
\newcommand{\bT}{\ensuremath{\mathbf{T}}}
\newcommand{\bU}{\ensuremath{\mathbf{U}}}
\newcommand{\bV}{\ensuremath{\mathbf{V}}}
\newcommand{\bW}{\ensuremath{\mathbf{W}}}
\newcommand{\bX}{\ensuremath{\mathbf{X}}}
\newcommand{\bY}{\ensuremath{\mathbf{Y}}}
\newcommand{\bZ}{\ensuremath{\mathbf{Z}}}
\newcommand{\bZero}{\ensuremath{\mathbf{0}}}
\newcommand{\bOne}{\ensuremath{\mathbf{1}}}
\newcommand{\bPsi}{\ensuremath{\mathbf{\Psi}}}
\newcommand{\bTheta}{\ensuremath{\mathbf{\Theta}}}
\newcommand{\bDelta}{\ensuremath{\mathbf{\Delta}}}
\newcommand{\bLambda}{\ensuremath{\mathbf{\Lambda}}}

%% -- Vectors with hat
\newcommand{\bahat}{\ensuremath{\hat{\bma}}}
\newcommand{\bbhat}{\ensuremath{\hat{\bmb}}}
\newcommand{\bchat}{\ensuremath{\hat{\bmc}}}
\newcommand{\bdhat}{\ensuremath{\hat{\bmd}}}
\newcommand{\behat}{\ensuremath{\hat{\bme}}}
\newcommand{\bfhat}{\ensuremath{\hat{\bmf}}}
\newcommand{\bghat}{\ensuremath{\hat{\bmg}}}
\newcommand{\bhhat}{\ensuremath{\hat{\bmh}}}
\newcommand{\bihat}{\ensuremath{\hat{\bmi}}}
\newcommand{\bjhat}{\ensuremath{\hat{\bmj}}}
\newcommand{\bkhat}{\ensuremath{\hat{\bmk}}}
\newcommand{\blhat}{\ensuremath{\hat{\bml}}}
\newcommand{\bmhat}{\ensuremath{\hat{\bmm}}}
\newcommand{\bnhat}{\ensuremath{\hat{\bmn}}}
\newcommand{\bohat}{\ensuremath{\hat{\bmo}}}
\newcommand{\bphat}{\ensuremath{\hat{\bmp}}}
\newcommand{\bqhat}{\ensuremath{\hat{\bmq}}}
\newcommand{\brhat}{\ensuremath{\hat{\bmr}}}
\newcommand{\bshat}{\ensuremath{\hat{\bms}}}
\newcommand{\bthat}{\ensuremath{\hat{\bmt}}}
\newcommand{\buhat}{\ensuremath{\hat{\bmu}}}
\newcommand{\bvhat}{\ensuremath{\hat{\bmv}}}
\newcommand{\bwhat}{\ensuremath{\hat{\bmw}}}
\newcommand{\bxhat}{\ensuremath{\hat{\bmx}}}
\newcommand{\byhat}{\ensuremath{\hat{\bmy}}}
\newcommand{\bzhat}{\ensuremath{\hat{\bmz}}}

%% -- Vectors with tilde
\newcommand{\batilde}{\ensuremath{\tilde{\bma}}}
\newcommand{\bbtilde}{\ensuremath{\tilde{\bmb}}}
\newcommand{\bctilde}{\ensuremath{\tilde{\bmc}}}
\newcommand{\bdtilde}{\ensuremath{\tilde{\bmd}}}
\newcommand{\betilde}{\ensuremath{\tilde{\bme}}}
\newcommand{\bftilde}{\ensuremath{\tilde{\bmf}}}
\newcommand{\bgtilde}{\ensuremath{\tilde{\bmg}}}
\newcommand{\bhtilde}{\ensuremath{\tilde{\bmh}}}
\newcommand{\bitilde}{\ensuremath{\tilde{\bmi}}}
\newcommand{\bjtilde}{\ensuremath{\tilde{\bmj}}}
\newcommand{\bktilde}{\ensuremath{\tilde{\bmk}}}
\newcommand{\bltilde}{\ensuremath{\tilde{\bml}}}
\newcommand{\bmtilde}{\ensuremath{\tilde{\bmm}}}
\newcommand{\bntilde}{\ensuremath{\tilde{\bmn}}}
\newcommand{\botilde}{\ensuremath{\tilde{\bmo}}}
\newcommand{\bptilde}{\ensuremath{\tilde{\bmp}}}
\newcommand{\bqtilde}{\ensuremath{\tilde{\bmq}}}
\newcommand{\brtilde}{\ensuremath{\tilde{\bmr}}}
\newcommand{\bstilde}{\ensuremath{\tilde{\bms}}}
\newcommand{\bttilde}{\ensuremath{\tilde{\bmt}}}
\newcommand{\butilde}{\ensuremath{\tilde{\bmu}}}
\newcommand{\bvtilde}{\ensuremath{\tilde{\bmv}}}
\newcommand{\bwtilde}{\ensuremath{\tilde{\bmw}}}
\newcommand{\bxtilde}{\ensuremath{\tilde{\bmx}}}
\newcommand{\bytilde}{\ensuremath{\tilde{\bmy}}}
\newcommand{\bztilde}{\ensuremath{\tilde{\bmz}}}

%% -- Lattices
\newcommand{\lLAMBDA}{\ensuremath{\mathbf{\Lambda}}}

%% -- MIMO specific macros
\newcommand{\MT}{\ensuremath{M_{T}}}
\newcommand{\MR}{\ensuremath{M_{R}}}
\newcommand{\SNR}{\ensuremath{\mathsf{SNR}}}
\newcommand{\SQNR}{\ensuremath{\mathrm{SQNR}}}
\newcommand{\SINR}{\ensuremath{\textsf{SINR}}}
\newcommand{\const}{\ensuremath{\setO}}
\newcommand{\mimoconst}{\ensuremath{\setO^{\MT}}}
\newcommand{\modorder}{\ensuremath{Q}}
\newcommand{\bHH}{\ensuremath{\bH^{H}}}
\newcommand{\bHbar}{\ensuremath{\bar{\bH}}}
\newcommand{\bQbar}{\ensuremath{\bar{\bQ}}}
\newcommand{\bRbar}{\ensuremath{\bar{\bR}}}
\newcommand{\bQa}{\ensuremath{\bQ_{\alpha}}}
\newcommand{\bRa}{\ensuremath{\bR_{\alpha}}}
\newcommand{\bIbar}{\ensuremath{\bar{\bI}}}
\newcommand{\bmybar}{\ensuremath{\bar{\bmy}}}
\newcommand{\bHi}{\ensuremath{\bH^{\dag}}}
\newcommand{\Pn}{\ensuremath{\sigma^{2}}}
\newcommand{\No}{\ensuremath{N_{o}}}
\newcommand{\VBLAST}{\ensuremath{\mbox{V-BLAST}}}
\newcommand{\Rate}{\ensuremath{R}}
\newcommand{\bsML}{\ensuremath{\bms^{\mathrm{ML}}}}
\newcommand{\bsZF}{\ensuremath{\bms^{\mathrm{ZF}}}}
\newcommand{\bsMMSE}{\ensuremath{\bms^{\mathrm{MMSE}}}}

%% -- VLSI things
\newcommand{\tclk}{\ensuremath{t_\mathrm{clk}}}
\newcommand{\fclk}{\ensuremath{f_\mathrm{clk}}}

%% -- Complexities
\newcommand{\CGr}{\ensuremath{C_{\mathrm{GR}}}}
\newcommand{\CVec}{\ensuremath{C_{\mathrm{Vec}}}}
\newcommand{\CRot}{\ensuremath{C_{\mathrm{Rot}}}}
\newcommand{\CMult}{\ensuremath{C_{\mathrm{Mult}}}}
\newcommand{\COMult}{\ensuremath{C_{\mathrm{OptMult}}}}
\newcommand{\CDiv}{\ensuremath{C_{\mathrm{Div}}}}
\newcommand{\CSqrt}{\ensuremath{C_{\sqrt{\cdot}}}}

%% -- NON-Math macros
\newcommand{\UMC}{$0.25\mu\mathrm{m}$ 1P/5M technology}

%% -- For NOTATION
\def\addnotation #1: #2{\parbox[t]{1.25cm}{$#1$\dotfill} \parbox[t]{3.50in}{#2} \vspace{0.0625cm} \\}
\def\addsymbol #1: #2{\parbox[t]{2cm}{$#1$} \parbox[t]{3in}{#2} \vspace{0.0625cm} \\}
\def\addacronym #1: #2{\parbox[t]{2cm}{\bf{#1}\dotfill} \parbox[t]{2.8in}{#2} \vspace{0.0625cm} \\}
   % include standard macros
\hyphenation{func-tion-al-ity}
\hyphenation{beam-forming}

%\newenvironment{code}{\begin{quote}}{\end{quote}}
%\newcommand{\sourcecode}[1]{\texttt{#1}}
\newcommand{\norm}{\xi}

\setcounter{tocdepth}{1} % Leveltiefe in Contents

% --- ansmath ------
\DeclareMathOperator\erfc{erfc}
\newcommand{\trace}[1]{\text{trace}\left(#1\right)}

%%%%%%%%%%%%%%%%%%%%%%%%%%%%%%%%%%%%%%%%%%%%%%%%%%%%%%%%%%%%%%%%%%%%%
%%%% BEGIN DOCUMENT %%%%%%%%%%%%%%%%%%%%%%%%%%%%%%%%%%%%%%%%%%%%%%%%%
%%%%%%%%%%%%%%%%%%%%%%%%%%%%%%%%%%%%%%%%%%%%%%%%%%%%%%%%%%%%%%%%%%%%%

\begin{document}
\frontmatter
\title{Distributed Multiuser Multihop Networks}
\author{Marc Müller}
\date{Spring 2014}

%=== TITLE =======================================================
\begin{titlepage}
  \thispagestyle{empty}

 \begin{figure}
 \subfigure{  \includegraphics[width=0.4\textwidth]{figures/ethlogo} } \hspace{1.5 cm}
 \subfigure{  \includegraphics[width=0.5\textwidth]{figures/narilogo} }
 \end{figure}

\vspace{3 cm}
%\vspace{\stretch{20}}

\noindent
\makebox[0pt][l]{
   \begin{minipage}{\textwidth}
      \flushright{
         \huge\bfseries
         \vspace{1.5 cm}
         Title of Master Thesis \\[1.2ex]
         \Large
         Subtitle of Master \\[1ex] Thesis
      }
      \noindent
      %\begin{center}
      \rule[-1ex]{\textwidth}{5pt}\\[2.5ex]      % draws black line 5pt high, \textwidth long, shifted 1ex down
      %\rule[-1ex]{\columnwidth}{5pt}\\[2.5ex]     
      %\end{center}
   \end{minipage}
}

\vspace{\stretch{7}}

 \begin{figure}[h]
 \begin{center}
 \subfigure{  \includegraphics[width=0.5\textwidth]{figures/fertigLuschtig} } 
 \end{center}
 \end{figure}

\vspace{\stretch{8}}

\noindent
\makebox[0pt][l]{
  \begin{minipage}{\textwidth}
    \flushright{
      \bfseries 
     \large Master Thesis \\
     \normalsize Summer 2009\\
     \Large
      	\vspace{1 cm}
	%\vspace{\stretch{5}}
	Muster Student\normalsize\\
      \vspace{0.5cm}
      Advisor:  Peter Schlaumeier\\ 
      Professor:  Prof. Dr. Armin Wittneben\\
      %Advisor:  J\"org Wagner\\ 
      [6ex]
	%\Large Master Thesis \normalsize \\
      %Summer 2009\\
    }
   \end{minipage}
}

\vfill

%\vspace{\stretch{1}}

\end{titlepage}

%=== Preamble ====================================================
\chapter*{Abstract}
\label{chap:abstract}
%----------------------------------------------------------
% Abstract
%----------------------------------------------------------
Here comes the abstract...

\chapter*{Preface}
\label{chap:preface}
%----------------------------------------------------------
% Preface
%----------------------------------------------------------

This Master Thesis is part of the graduate study at the Department
of Information \linebreak
Technology and Electrical Engineering (D-ITET) at the
Swiss Federal Institute of Technology (ETH) Zurich.
\\
\par The author certifies that this Master Thesis, and the research
to which it refers, are the product of the author's own work, and that
any ideas or quotations from the work of other people, published or
otherwise, are fully acknowledged in accordance with the standard
referring practices of the discipline.

\vspace{12mm}
\noindent
Marc Müller
\vspace{15mm}
\begin{table}[h!]
\begin{tabular}{p{55mm}}
\hline \\
\end{tabular}
\end{table}

%----------------------------------------------------------
% Acknowledgements
%----------------------------------------------------------

\chapter*{Acknowledgements}
\label{chap:acknowledgements}

I would like to thank blablabla

\cleardoublepage

%----------------------------------------------------------
% People
%----------------------------------------------------------
\begin{center}
\begin{tabular}{l@{ }l@{\hspace{1cm}}l}
\emph{Author:} & Marc Müller & marcmue@student.ee.ethz.ch\\
\emph{Advisor:}& Yahia Hassan & yehia.hassan@nari.ee.ethz.ch\\
\emph{Professor:}& Prof. Dr. Armin Wittneben & wittneben@nari.ee.ethz.ch\\
\end{tabular}
\end{center}

\pagebreak
 % Abstract, Preface, Acknowledgements

%=== Verzeichnisse ===============================================
\tableofcontents
\listoffigures 
\listoftables
\listofalgorithms

%=== Chapters ====================================================
\mainmatter
\renewcommand{\baselinestretch}{1}\normalsize % Zeilenabstand
%\settocdepth{section} % Kapitel Sectionweise in Contents
\chapter{Information Theory}


\section{general Channel}
//by book of Tse\\
We use $x(t)$ as the transmitted signal and $y(t)$ as the received signal. Since the signal is affected by different objects in its path the received signal can be written as a function of the resulting different paths
\begin{equation}
y(t)=\sum_i{a_i(f,t)x(t-\tau_i(t))}.
\end{equation}
$a_i(f,t)$ describes the attenuation of each path which is mostly dependent on frequency and time. we can omit these dependencies if we can assume sufficient narrow bandwidth for a frequency flat channel and with relatively stationary objects a time flat channel, objects are transmitters, receivers, scatterers and so on. $\tau_i(t)$ is the delay in time of the path and is constant in a time flat channel.\\
To distinguish between these different channels we introduce some measures:
\begin{description}
	\item[Delay Spread $T_d$] describes the time difference from the first to the last arrival of the signal at the receiver.
	\item[Doppler Spread $D_s$] describes the maximum frequency difference which is introduced by Doppler shifts.
	\item[coherence Bandwidth $W_c$] is the inverse of the delay spread and a measure of the frequency flatness of the channel. It is the bandwidth at which we can assume that the channel is frequency independent.
	\item[coherence Time $T_c$] is the inverse of the Doppler spread and a measure for the time flatness of the channel e.g. it is the time over which the channel is time invariant. 
\end{description}
// image

\subsection{The equivalent baseband representation}
If we have a frequency-flat and time-flat channel we can represent our channel with a time varying channel gain $h$ which leads to 
\begin{equation}
	y(t) = h \cdot x(t)
\end{equation} for our input output relation. But it is often simpler to use a complex baseband model. if we sample the system at a sufficiently high rate we get 
\begin{equation}
	y[m] = \sum_l{h_l \cdot x[m]}
\end{equation}
where $m$ describes discrete time $x[m]$ and $y[m]$ the transmitted respectively the received signal.  

\subsection{Channel Capacity}
the channel capacity is a good figure of merit to describe the optimal performance of a communication system. It gives a maximum data rate at which we could transmit error-free when using appropriate coding.
\begin{description}
	\item[AWGN-channel:] in an AWGN channel we experience no fading and the received signal consists just the sent signal with additive white Gaussian noise $y=x+n$. The capacity depends only on the transmit $P$ versus noise Power $N_0$:
	\begin{equation}
	C=log_2\Biggr(1+\frac{P}{N_0}\Biggl) \text{     bits/s/Hz}
	\end{equation}
	\item[time- and frequency-flat fading channel:] in a one tab fading channel the capacity is a function of the tab gain $h$:
	\begin{equation}C=log_2\Biggr(1+|h|^2\frac{P}{N_0}\Biggl) \text{     bits/s/Hz}
	\end{equation}
\end{description}

\section{MIMO}
IN the thrive for higher data rates the limits of the basic single Antenna communication channels became obvious. The bandwidth limitation due to scarcity and cost and power constraints due to physics and regulations called for different approaches. MIMO uses the spatial dimension of the channel to further increase reliability and speed.

\subsection{from SISO to MIMO}
We are still using the assumptions of frequency- and time-flat channels. In the SISO case the discrete complex baseband channel tab $h$ characterized the channel sufficiently. If we now use multiple antennas an the transmitter and receiver side of the channel and make sure that they are sufficiently spaced it establishes multiple distinct channels due to small-scale fading processes. In this case our discrete complex baseband channel tab $h$ is no longer a single coefficient but a matrix $\mathbf{H}\in\mathbb{C}^{N\times M}$ with $M$ number of transmit antennas and $N$ number of receive antennas. 
\begin{equation}
	\mathbf{H} = \begin{bmatrix}
		h_{11} & h_{12} & \hdots & h_{1M} \\
		h_{21} & h_{22} & \hdots & h_{2M} \\
		\vdots & \vdots & \hdots & \vdots \\
		h_{N1} & h_{N2} & \hdots & h_{NM} \\
	\end{bmatrix}
\end{equation}

\\
//figure \\
the input output relation of a MIMO channel is then written as 
\begin{equation}
\mathbf{y}[m] = \mathbf{Hx}[m] + \mathbf{w}[m]
\end{equation}
where $\mathbf{x}\in\mathbb{C}^M , \mathbf{y}\in\mathbb{C}^N$ and $\mathbf{w}\sim\mathcal{CN}(0,N_0\mathbf{I}_N)$


\subsection{MIMO Capacity}
the capacity of a MIMO channel without any precodeing can be written as:
\begin{equation}
	C = \log_2{\det{\Biggr[\mathbf{I}_N + \frac{\text{SNR}}{M}\mathbf{HH}^H\Biggl]}}
\end{equation}
with $\text{SNR}$ as total transmit power $P$ over noise power $N_0$.


% For a SISO channel in a Rayleigh fading channel we get a bit error rate off
% \begin{equation}
% 	P_b(|h|^2) = 1/2 \erfc{\Biggl(\sqrt{|h|^2\frac{\bar{E_b}}{N_0}}\Biggr)}.
% \end{equation}
% , if we now send the same signal over $N$ statistically independent channels we get a BER of
% \begin{equation}
% 	P_b(|h|^2) = 1/2 \erfc{\Biggl(\sqrt{\sum_{j=1}^N{|h_j|^2\frac{E_s}{N_0}}}\Biggr)}
% \end{equation}
% for the equivalent SISO channel.



\chapter{Optimization}

The question now is, is it possible to increase the Throughput of this channel with knowledge of the channel Matrix $H$. and to investigate this we use following Channel Representation and we are fully aware of its implications and imperfections.
\\// description of our channel

\section{Waterfilling}
If full CSI and full inter antenna coordination e.g. power and data can be freely distributed over all transmit antennas there is a single easy to calculate optimal solution; water-filling. Water filling is designed to diagonalize the channel with help of the SVD. one huge drawback to it is that we need the antenna coordination on the transmitter side, so it is only hardly feasible for a multiuser case.\\
As mentioned water-filling starts with the SVD of the channel $\mathbf{H} = \mathbf{USV}^H$. It then uses $\mathbf{\tilde{x}} = \mathbf{Vx}$ and $\mathbf{\tilde{y}} = \mathbf{U}^H\mathbf{y}$ to rotate the transmit and receive vectors. $\mathbf{\tilde{w}} = \mathbf{U}^H\mathbf{w}$ is the rotate noise vector.
\begin{equation}
	\mathbf{\tilde{y}} 
	= \mathbf{U}^H\mathbf{USV}^H\mathbf{Vx}+\mathbf{U}^H\mathbf{w}}
	= \mathbf{Sx + \tilde{w}}.
\end{equation}
Now we have multiple parallel SISO "eigen-subchannel" systems. The waterfilling now says how to distribute the transmit power $P_{tot}$ over the different subchannels. The idea is to give the "strongest" channel the most power and the "weakest" the least amount.\\

// picture

The question to answer how to distribute the power to maximize throughput $C$:
\begin{equation}
	\begin{aligned}
		& \underset{\mathbf{P}}{\text{maximize}}
		& & \sum_k{\log_2{\Biggr(1+\frac{P_k|\tilde{h}_k|}{N_0}\Biggl)}} \\
		& \text{subject to}
		& & P_k\geq 0,\,\forall P_k\in \diag{\mathbf{P}}, \\
		&&& \trace{\mathbf{P}}\leq P_{tot}.
	\end{aligned}
\end{equation}
with noise Power $N_0$ and $\tilde{h}_k$ the complex channel gain of subchannel $k$
We define $(\cdot)^+ := \max{\cdot,0}$ the power allocation
\begin{equation}
	P_k^* = \Biggr(\frac{1}{\lambda}-\frac{N_0}{|\tilde{h}_k|^2}\Biggl)^+
\end{equation}
with the water-level $\lambda$ chosen such that the power constraint is met
\begin{equation}
	\frac{1}{M}\sum_{k=1}^M{\Biggr(\frac{1}{\lambda}-\frac{N_0}{|\tilde{h}_k|^2}\Biggl)^+}
\end{equation}

\section{Multiuser Case}
If antenna coordination is not possible for example in a multiuser environment the water filling solution is not reachable. Jindal proposes an algorithm his calling iterative waterfilling. The goal of this algorithm is to calculate the effective noise (Gaussian noise plus interference) for each user and then use water filling over these adapted channels. But since changing the power of one user changes the effective noise of all other users its mandatory to iterate until the algorithm converges. His algorithm optimizes with respect to an SIC receiver.
The papers goal is to solve the problem of beam-forming in a broadcast(BC) or downlink channel, so exactly the opposite direction than we are discussing in this report. But as //paper describes there is a duality of these two channels and we are able to use the optimized results of the multiple access channel (MAC) or uplink-channel to optimize the BC. His work is using this duality to solve the easier problem of the MAC and then transform it back to the BC problem, ant thus solving the throughput maximization problem for the multiuser case.

\begin{enumerate}
	\item[init] with $\mathbf{Q^{(0}}=\mathbf{0}$
	\item Generate effective channels for each user
		\begin{equation}
			\mathbf{G}_k^{(n)}=\mathbf{H}_k\Biggr(\mathbf{I}+\sum_{j\neq k}{\mathbf{H}_j^H\mathbf{Q}_k^{(n-1)}\mathbf{H}_j}\Biggl)^{-1/2}
		\end{equation}
	\item Treating these effective channels as parallel, noniterfering channels and obtain covariance matrices $\mathbf{S}_k^{(n)}$ by waterfilling with total power $P_{ot}$
		\begin{equation}
			\begin{aligned}
				& \underset{\mathbf{\mathbf{S}_k}}{\text{maximize}}
				& & \sum_k{\log_2{\Biggr(\mathbf{I}+\Bigr(\mathbf{G}_k^{(n)}\Bigl)^H\mathbf{S}_k\mathbf{G}_k^{(n)}\Biggl)}} \\
				& \text{subject to}
				& & \mathbf{S}_k \geq 0,\\
				&&& \sum_k{\trace{\mathbf{S}_k}}\leq P_{tot}.
			\end{aligned}
		\end{equation}
	\item Compute the updated covariance matrices $\mathbf{Q}_k^{(n)}$ as
		\begin{equation}
			\mathbf{Q}_k^{(n)} = \frac{1}{M}\mathbf{S}_k^{(n)}+\frac{M-1}{M}\mathbf{Q}_k^{(n-1)}
		\end{equation}
	where $M$ is the number of users.
\end{enumerate}


// figure (compare wf to iwf to equal power allocation)
\section{Optimizing for MMSE with $K$ users}

We have a channel with $K$ users which all have one single antenna and one receiver with $M$ antennas. Our system equation then is 
\begin{equation}
	\mathbf{y}^{(M\times 1)} = \mathbf{H}^{(M\times K)}\cdot \mathbf{x}^{(K\times 1)}.
\end{equation}
As discussed before the receive covariance matrix is 
\begin{equation}
	\mathbf{K}_e = \Biggr(\mathbf{\sqrt{P}H}^H\frac{1}{\sigma^2}\mathbf{H\sqrt{P}}+\mathbf{I}\Biggl)^{-1}
\end{equation}
with $\mathbf{P} = \diag{P_1,\dots,P_K}$ the power loading matrix where $P_k$ is the power of user $k$.\\
We then calculate the signal to noise and interference ratio (SINR) and the rate $R$ of each user as
\begin{align}
	\text{SINR}_k = \frac{1}{(\mathvf{K}_e)_{k,k}}\\
	R_k = \log_2{1+\text{SINR}_k}.
\end{align}
We further define a fixed power constraint $P_{tot}$ so that $ \sum_k{P_k} \leq P_{tgt}$ and $P_k\leq 0\,,\forall k$. 

\subsection{Gradient search}

The idea of the steepest descent algorithm is to find the global maximum or minimum of a function by searching along the gradient of the function.
\begin{algorithm}
	\KwIn{$\mathbf{X^{(0)}}$, $f(\mathbf{X})$, $\nabla f(\mathbf{X})$, $\mu$(step size)}
	\KwOut{$\mathbf{X}^{(\text{end})}$}
	\Repeat{$\nabla f(\mathbf{X}^{(n)}) = 0$}{
		$
			\mathbf{X}^{(n+1)}=\mathbf{X}^{(n)}-\mu\nabla f(\mathbf{X}^{(n)})
		$\;
	}		
\end{algorithm}\\
To optimize for the rates of a given channel $\mathbf{H}$ we define
\begin{equation}
	f_p(\mathbf{P})=\bigr\Vert{(R(\mathbf{P}))\bigl\Vert}_p = \Biggr(\sum_k{(R_k(\mathbf{P}))^p}\Biggl)^{1/p}.
\end{equation}
as the $p$-norm of the user rates. and calculate its gradient $\nabla f_p(\mathvf{P})$ as
\begin{equation}
	\nabla f_p(\mathbf{P}) =
	\begin{bmatrix}
		\frac{\partial f_p(\mathbf{P})}{\partial P_1}\\
		\frac{\partial f_p(\mathbf{P})}{\partial P_2}\\
		\vdots\\
		\frac{\partial f_p(\mathbf{P})}{\partial P_K}.
	\end{bmatrix}
\end{equation}
This gives us an iterative search algorithm
\begin{equation}
	\mathbf{P}^{n+1} = \mathbf{P}^n + \mu \nabla \bigr\Vert{(R(\mathbf{P}))\bigl\Vert}_p
\end{equation}
where $(\cdot)^n$ denotes the $n$-th iteration of the search algorithm and $\mu$ an arbitrary step size.
In the following we use $\partial_{X_j}(f_p) = \frac{\partial f_p(\mathbf{X})}{\partial X_j}$ and $R = R(\mathbf{X})$.

We now calculate the gradient of
\begin{equation}
	\bigr\Vert{(R)\bigl\Vert}_p = \Biggr(\sum_k{(R_k)^p}\Biggl)^{1/p}.
\end{equation}
as follows
\begin{align}
	\nabla \bigr\Vert{(R(\mathbf{P}))\bigl\Vert}_p &= \nabla \Biggr[ \Biggr(\sum_k{(R_k(\mathbf{P}))^p}\Biggl)^{1/p} \Biggl]\\
	&= \begin{bmatrix}
		\partial_{p_1}\\\partial_{p_2}\\\vdots\\\partial_{p_M}
	\end{bmatrix}
	\cdot \Biggr(\sum_k{(R_k(\mathbf{P}))^p}\Biggl)^{1/p}.
\end{align}

Due to the total power constraint $P_{tot}$ and that no user can send with negative power we have a bounded optimization problem. So we need to convert it into an unbounded one:
\begin{enumerate}
 	\item to guarantee $P_k \geq 0 \, \forall k\in\diag{\mathbf{P}}$ we can introduce a diagonal matrix $\mathbf{X}$ as $X_k^2 = P_k$ with $X_k$ is the $k$-th diagonal element of $\mathbf{X}$.
 	\item to guarantee $\trace{\mathbf{P}}\leq P_{tot}$ we tighten the condition to $\trace{\mathbf{P}}= P_{tot}$ and calculate the rate with:
 	\begin{align}
	R_k &= -\log_2\Biggr(\mathbf{\Phi}_{k,k}^{-1}\Biggl)\\
	&= -\log_2\Biggr(\biggr[ \sqrt{\mathbf{P}}\mathbf{H}^H\frac{1}{\sigma^2}\mathbf{H}\sqrt{\mathbf{P}} + \mathbf{I_M} \biggl]_{k,k}^{-1}\Biggl)\\
	&= -\log_2\Biggr(\biggr[ \sqrt{\mathbf{P} \frac{P_{tot}}{\trace{\mathbf{P}}}}\mathbf{H}^H\frac{1}{\sigma^2}\mathbf{H} \sqrt{\mathbf{P} \frac{P_{tot}}{\trace{\mathbf{P}}}} + \mathbf{I_M} \biggl]_{k,k}^{-1}\Biggl)\\
	&= -\log_2\Biggr(\biggr[ \mathbf{X} \sqrt{\frac{P_{tot}}{\trace{\mathbf{X}^2}}} \mathbf{H}^H\frac{1}{\sigma^2}\mathbf{H} \mathbf{X} \sqrt{\frac{P_{tot}}{\trace{\mathbf{X}^2}}} + \mathbf{I_M} \biggl]_{k,k}^{-1}\Biggl)\\
	&= -\log_2\Biggr(\biggr[\frac{P_{tot}}{\trace{\mathbf{X}^2}} \mathbf{X}\mathbf{H}^H\frac{1}{\sigma^2}\mathbf{H} \mathbf{X} + \mathbf{I_M} \biggl]_{k,k}^{-1}\Biggl)
\end{align}
\end{enumerate}

We now can calculate the partial derivatives of the Rates:
\begin{align}
	\partial_{x_j} \bigr\Vert{(R(\mathbf{P}))\bigl\Vert}_p
	&=\partial_{x_j}\Biggr[\Biggr(\sum_k{(R_k(\mathbf{P}))^p}\Biggl)^{1/p}\Biggl]\\
	&=\frac{1}{p}\Biggr(\sum_k{(R_k(\mathbf{P}))^p}\Biggl)^{1/p-1} \cdot \sum_k{\Biggr[p\cdot(R_k(\mathbf{P}))^{p-1} \cdot \partial_{x_j}(R_k(\mathbf{P}))\Biggl]}.
\end{align}
with the partial derivative of the $k$-th users rate
\begin{align}
	\partial_{x_j}(R_k) &= \partial_{x_j}\bigr(-\log_2{\mathbf{\Phi}_{k.k}^{-1}}\bigl)\\
	&=\frac{-1}{\log{2}}\frac{1}{\mathbf{\Phi}_{k,k}^{-1}} \cdot \partial_{x_j}\bigr(\mathbf{\Phi}_{k,k}^{-1}\bigl).
\end{align}
We use
\begin{equation}
	\partial_{x_j}\bigr(\mathbf{\Phi}_{k,k}^{-1}\bigl) = \Biggr(-\mathbf{\Phi}^{-1} \cdot \partial_{x_j}\bigr(\mathbf{\Phi}^{-1}\bigl) \cdot \mathbf{\Phi}^{-1}\Biggl)_{k,k}.
\end{equation} to calculate the partial derivative of an inverse matrix, such that the derivative of the of the inverse matrix $\mathbf{\Phi}$ becomes
\begin{align}
	\partial_{x_j}\bigr(\mathbf{\Phi}^{-1}\bigl) &= \partial_{x_j} \biggr(\frac{P_{tot}}{\trace{\mathbf{X}^2}} \mathbf{X}\mathbf{H}^H\frac{1}{\sigma^2}\mathbf{H} \mathbf{X} + \mathbf{I_M} \biggl)\\
	&= \partial_{x_j} \biggr(\frac{P_{tot}}{\trace{\mathbf{X}^2}} \mathbf{X}\mathbf{H}^H\frac{1}{\sigma^2}\mathbf{H} \mathbf{X}\biggl)\\
	&= P_{tot} \cdot \partial_{x_j} \biggr(\frac{ \mathbf{X}\mathbf{H}^H\frac{1}{\sigma^2}\mathbf{H} \mathbf{X} }{\trace{\mathbf{X}^2}}\biggl)\\
	&= P_{tot} \cdot \Biggr(\frac{\partial_{x_j}\bigr(\mathbf{X}\mathbf{H}^H\frac{1}{\sigma^2}\mathbf{H}\mathbf{X}\bigl)}{\trace{\mathbf{X}^2}} - \frac{\mathbf{X}\mathbf{H}^H\frac{1}{\sigma^2}\mathbf{H}\mathbf{X} \cdot \partial_{x_j}\bigr(\trace{\mathbf{X}^2}\bigl)}{\bigr(\trace{\mathbf{X}^2}\bigl)^2}\Biggl).
\end{align}
The partial derivative of a matrix $\frac{\partial(\mathbf{X})}{\partial X_{i,j}} = E_{i,j}$ with $E_{i,j}$ the single entry matrix, 1 at $(i,j)$ and zero elsewhere. Ant thus
\begin{align}
	\partial_{x_j}\bigr(\mathbf{X}\mathbf{H}^H\frac{1}{\sigma^2}\mathbf{H}\mathbf{X}\bigl) &= \mathbf{E}_{j,j}\mathbf{H}^H\frac{1}{\sigma^2}\mathbf{HX} + \mathbf{XH}^H\frac{1}{\sigma^2}\mathbf{HE}_{j,j}\\
	\partial_{x_j}\bigr(\trace{\mathbf{X}^2}\bigl) &= 2\mathbf{XE}_{j,j}.
\end{align}
Finally we combine it all together for the $j$-th partial derivative of the $p$-norm of the user rates:
\begin{multline}
	\partial_{x_j} \bigr\Vert{(R)\bigl\Vert}_p =\\
	\Biggr(\sum_k{(R_k)^p}\Biggl)^{1/p-1} \cdot \sum_k{\Biggr[\frac{1}{\log{2}}\frac{(R_k)^{p-1}}{\mathbf{\Phi}_{k,k}^{-1}} \cdot \Biggr(\mathbf{\Phi}^{-1} \cdot \\ P_{tot} \cdot \Biggr(\frac{\mathbf{E}_{j,j}\mathbf{H}^H\frac{1}{\sigma^2}\mathbf{HX} + \mathbf{XH}^H\frac{1}{\sigma^2}\mathbf{HE}_{j,j}}{\trace{\mathbf{X}^2}} -\\ \frac{\mathbf{X}\mathbf{H}^H\frac{1}{\sigma^2}\mathbf{H}\mathbf{X} \cdot 2\mathbf{XE}_{j,j}}{\bigr(\trace{\mathbf{X}^2}\bigl)^2}\Biggl)
	\cdot \mathbf{\Phi}^{-1}\Biggl)_{k,k}\Biggl]}
	%\partial_{x_j}(R_k)) = \frac{P_{tot}}{\log{2}\cdot\mathbf{\Phi}_{k,k}^{-1}} \Biggr(\mathbf{\Phi}^{-1} \cdot \Biggr(\frac{\mathbf{E}_{j,j}\mathbf{H}^H\frac{1}{\sigma^2}\mathbf{HX} + \mathbf{XH}^H\frac{1}{\sigma^2}\mathbf{HE}_{j,j}}{\trace{\mathbf{X}^2}} - \frac{\mathbf{X}\mathbf{H}^H\frac{1}{\sigma^2}\mathbf{H}\mathbf{X} \cdot 2\mathbf{XE}_{j,j}}{\bigr(\trace{\mathbf{X}^2}\bigl)^2}\Biggl) \cdot \mathbf{\Phi}^{-1}\Biggl)_{k,k}
\end{multline}
% and finally for maximizing the sum Rate:
% \begin{equation}
% 	\nabla \bigr\Vert{(R(\mathbf{P}))\bigl\Vert}_1 =
% 	\begin{bmatrix}
% 		\sum_k{\Biggr[
% 			\frac{-1}{\log{2}}\frac{1}{\mathbf{\Phi}_{k,k}^{-1}} \cdot P \cdot \Biggr(\frac{\mathbf{E}_{1,1}\mathbf{H}^H\frac{1}{\sigma^2}\mathbf{HX} + \mathbf{XH}^H\frac{1}{\sigma^2}\mathbf{HE}_{1,1}}{\trace{\mathbf{X}^2}} - \frac{\mathbf{X}\mathbf{H}^H\frac{1}{\sigma^2}\mathbf{H}\mathbf{X} \cdot 2\mathbf{XE}_{1,1}}{\bigr(\trace{\mathbf{X}^2}\bigl)^2}\Biggl)
% 		\Biggl]}\\
% 		\sum_k{\Biggr[
% 			\frac{-1}{\log{2}}\frac{1}{\mathbf{\Phi}_{k,k}^{-1}} \cdot P \cdot \Biggr(\frac{\mathbf{E}_{2,2}\mathbf{H}^H\frac{1}{\sigma^2}\mathbf{HX} + \mathbf{XH}^H\frac{1}{\sigma^2}\mathbf{HE}_{2,2}}{\trace{\mathbf{X}^2}} - \frac{\mathbf{X}\mathbf{H}^H\frac{1}{\sigma^2}\mathbf{H}\mathbf{X} \cdot 2\mathbf{XE}_{2,2}}{\bigr(\trace{\mathbf{X}^2}\bigl)^2}\Biggl)
% 		\Biggl]}\\
% 		\vdots\\
% 		\sum_k{\Biggr[
% 			\frac{-1}{\log{2}}\frac{1}{\mathbf{\Phi}_{k,k}^{-1}} \cdot P \cdot \Biggr(\frac{\mathbf{E}_{m,m}\mathbf{H}^H\frac{1}{\sigma^2}\mathbf{HX} + \mathbf{XH}^H\frac{1}{\sigma^2}\mathbf{HE}_{m,m}}{\trace{\mathbf{X}^2}} - \frac{\mathbf{X}\mathbf{H}^H\frac{1}{\sigma^2}\mathbf{H}\mathbf{X} \cdot 2\mathbf{XE}_{m,m}}{\bigr(\trace{\mathbf{X}^2}\bigl)^2}\Biggl)
% 		\Biggl]}
% 	\end{bmatrix}
% \end{equation}

\subsection{maximum throughput}
We want to maximize the sum-rate for a given channel $\mathbf{H}$ and a fixed power constraint $P_{tot}$:
\begin{equation}
	\begin{aligned}
		& \underset{\mathbf{P}}{\text{maximize}}
		& & \sum_k{R_k(\mathbf{P})} \\
		& \text{subject to}
		& & P_k\geq 0,\,\forall P_k\in \diag{\mathbf{P}}, \\
		&&& \trace{\mathbf{P}}\leq P.
	\end{aligned}
\end{equation}
this optimizes the diagonal precoding matrix $\mathbf{P}$ with $\mathbf{P} = \diag{P_1,\dots,P_K},\,\forall P_k \geq 0$. Using the previously calculated gradient of $\bigr\Vert{(R(\mathbf{P}))\bigl\Vert}_p$ with $p = 1$ for $\bigr\Vert{(R(\mathbf{P}))\bigl\Vert}_1 = $

As sated before the $p$-norm gives us the possibility to use the same gradient for different applications; setting $p=1$ gives us $\bigr\Vert{(R(\mathbf{P}))\bigl\Vert}_1 = \sum_k{R_k}$ to optimize for the sum-rate:
\begin{multline}
	\partial_{x_j} \Bigr(\sum_k{R_k}\Bigl) =
	\sum_k{\Biggr[\frac{1}{\log{2}}\frac{1}{\mathbf{\Phi}_{k,k}^{-1}} \cdot \Biggr(\mathbf{\Phi}^{-1} \cdot P_{tot} \\\cdot \Biggr(\frac{\mathbf{E}_{j,j}\mathbf{H}^H\frac{1}{\sigma^2}\mathbf{HX} + \mathbf{XH}^H\frac{1}{\sigma^2}\mathbf{HE}_{j,j}}{\trace{\mathbf{X}^2}} - \frac{\mathbf{X}\mathbf{H}^H\frac{1}{\sigma^2}\mathbf{H}\mathbf{X} \cdot 2\mathbf{XE}_{j,j}}{\bigr(\trace{\mathbf{X}^2}\bigl)^2}\Biggl)
	\cdot \mathbf{\Phi}^{-1}\Biggl)_{k,k}\Biggl]}
\end{multline}

\subsection{maximum min-rate}
To maximize the min-rate we would use $p=-\infty$ for the$\Vert\cdot\Vert_{min}$ if there is no per user power constraint this gives us equal throughput for each user. But the min-norm is non differentiable and thus unusable for our gradient search algorithm. In search of a feasible solution we use the $-30$-norm, this gives us a good trade-of between stability and preciseness of the results. a $p$-Value of $-100$ could work to but the risk of getting non differentiable is too high.
We thus optimize
\begin{equation}
	\begin{aligned}
		& \underset{\mathbf{P}}{\text{maximize}}
		& & \Bigr\Vert{R(\mathbf{P})}\Bigl\Vert_{-30} \\
		& \text{subject to}
		& & P_k\geq 0,\,\forall P_k\in \diag{\mathbf{P}}, \\
		&&& \trace{\mathbf{P}}\leq P.
	\end{aligned}
\end{equation}
which is the maxmin algorithm.

\subsection{optimizing the gradient algorithm}
One problem that occurred during the simulations was that at some point the gradient search got unstable, in our case it started to fluctuate and stopped moving towards the optimum. The reason of this bad and unwanted behavior lies in the step size of the function, and so there are two ways to fix it. The first and easy way is just to reduce the step size and increase the number of iterations accordingly, this way the problem occurs later and weaker and has less impact on the end result. The second much more complicated solution is to introduce adaptive step size.

Adaptive step size means, that for each iteration we search the optimal step size which allows us to take the biggest step towards our goal. So it should not be possible to choose a step size which gives us a worse result than we had on the iteration before.
\begin{algorithm}
	\KwIn{$\mathbf{X}^{(0)}$, $f(\mathbf{X}^{(0)})$, $\nabla f(\mathbf{X}^{(0)})$, $\mu^{(0)}$, $\tau$ (step size multiplicator)}
	\eIf{$f(\mathbf{X}^{(n)}) \leq f(\mathbf{X}^{(n+1)})$}{
		\Repeat{$f(\mathbf{X}^{(n)}) \leq f(\mathbf{X}^{(n-1)})$}{
			$\mu^{(n)} = \mu^{(n-1)} \cdot \tau$\;
			$\mathbf{X}^{(n)}=\mathbf{X}^{(0)}-\mu^{(n)}\nabla f(\mathbf{X}^{(0)})$\;
			calculate: $f(\mathbf{X}^{(n)})$\;
		}
		\Return{$\mathbf{X}^{(n-1)}$}
	}
	{
		\Repeat{$f(\mathbf{X}^{(n)}) > f(\mathbf{X}^{(n-1)})$}{
		$\mu^{(n)} = \mu^{(n-1)} \cdot 1/\tau$\;
		$\mathbf{X}^{(n)}=\mathbf{X}^{(0)}-\mu^{(n)}\nabla f(\mathbf{X}^{(0)})$\;
		calculate: $f(\mathbf{X}^{(n)})$\;
		}
		\Return{$\mathbf{X}^{(n)}$}
	}
\end{algorithm}
With this Algorithm we can prevent any fluctuation in the gradient and force it to go the fastest possible way. We do not have to calculate the gradient as often and can thus reduce the needed calculation time. Another nice side effect is, that we implicitly know that we reached the maximum or minimum when the step size is decreased to zero without finding any better points than the previous one.
\section{optimizing VBlast}

\subsection{numeric Gradient}
\subsubsection{max sumRate}
\subsubsection{max minRate}


\section{optimizing for SINR target}
\subsection{fodor}

\subsection{The Algorithm}
\begin{enumerate}
	\item calculate the effective interference for MMSE processing:
		\begin{equation}
			\begin{aligned}
				\zeta_{k,s} = \Biggl\{\Biggl(d^{-\rho}_{k,k}\chi_{k,k}H^H_{k,k}\biggl(\sum_{j\neq k}{P_jd^{-\rho}_{k,j}\chi_{k,j}H_{k,j}T_jT_j^HH^H_{k,j}}\\
				+N_t\sigma^2_nI\biggr)^{-1}H_{k,k}+\frac{1}{P_k}I\Biggr)^{-1}\Biggr\}^{(s,s)}
			\end{aligned}	
		\end{equation}

	\item calculate the optimal loading matrix
		\begin{equation}
			(T_k)^{(s,s)} = \sqrt{\frac{\zeta_{k,s}N_t}{\sum_{j=1}^{N_t}\zeta_{k,j}}}\;\forall s\in[1,N_t]
		\end{equation}

	\item calculate used Power
		\begin{equation}
			P_k = \frac{\zeta_{k,s}}{\vert(T_k)^{(s,s)}\vert^2}(\gamma_{tgt}+1)\;\forall k
		\end{equation}

	\item[n.] until no more change

\end{enumerate}

\subsection{analytical gradient}
We calculate the power minimization on the basis of the analytical gradient for throughput maximization.
\begin{equation}
	\mathbf{P} = \arg{\min_{\begin{matrix}
		P_k\geq 0\,\forall P_k\in \diag{\mathbf{P}}\\
		R_k = R_{tgt}\,\forall k
		\end{matrix}}{\sum_k{P_k}}}
\end{equation}The idea here was to minimize the difference between the target Rates and the calculated rate for the channel, only by changing the diagonal Values of $\mathbf{P}$.
So our problem gets reformed into
\begin{equation}
	\mathbf{P} = \arg{\max_{P_k\geq 0\,\forall P_k\in \diag{\mathbf{P}}}{\sum_k{\bigr(R_k-R_{tgt}\bigl)}}}
\end{equation}
Again we use $\mathbf{X}^2 = \mathbf{P}$ to make sure to only get positive power Values but since we do not care about some power power budget $P_{tot}$ we can removed the condition $\trace{\mathbf{P}}\leq P_{tot}$. $\mathbf{\Phi}$ then becomes
\begin{equation}
	\mathbf{\Phi} = \mathbf{X}\mathbf{H}^H\frac{1}{\sigma^2}\mathbf{H} \mathbf{X} + \mathbf{I_M}
\end{equation}
and the for the Rate of user $k$
\begin{equation}
	R_k = -\log_2\Biggr(\mathbf{\Phi}_{k,k}^{-1}\Biggl) = -\log_2\Biggr(\biggr[\mathbf{X}\mathbf{H}^H\frac{1}{\sigma^2}\mathbf{H} \mathbf{X} + \mathbf{I_M} \biggl]_{k,k}^{-1}\Biggl)
\end{equation}
With $p=2$ we calculate the distance such that the sign does not mess up the gradient and the norm to be able to reuse the code from the throughput maximization algorithm, we get:
\begin{equation}
	\nabla \bigr\Vert{(R(\mathbf{P}))\bigl\Vert}_2 = \nabla \Biggr[ \Biggr(\sum_k{(R_k-R_{tgt})^2}\Biggl)^{1/2} \Biggl]
\end{equation}
and for the partial derivative
\begin{align}
	\partial_{x_j} \bigr\Vert{(R(\mathbf{P}))\bigl\Vert}_2
	&=\partial_{x_j}\Biggr[\Biggr(\sum_k{(R_k-R_{tgt})^2}\Biggl)^{1/2}\Biggl]\\
	&=\frac{1}{2}\Biggr(\sum_k{(R_k-R_{tgt})^2}\Biggl)^{-1/2} \cdot \sum_k{\Biggr[2(R_k-R_{tgt}) \cdot \partial_{x_j}(R_k-R_{tgt})\Biggl]}.
\end{align}
where
\begin{align}
	\partial_{x_j}(R_k-R_{tgt}) &= \partial_{x_j}(R_k)\\
	&= \partial_{x_j}\bigr(-\log_2{\mathbf{\Phi}_{k.k}^{-1}}\bigl)\\
	&=\frac{-1}{\log{2}}\frac{1}{\mathbf{\Phi}_{k,k}^{-1}} \cdot \partial_{x_j}\bigr(\mathbf{\Phi}_{k,k}^{-1}\bigl)
\end{align}
\begin{equation}
	\partial_{x_j}\bigr(\mathbf{\Phi}_{k,k}^{-1}\bigl) = \Biggr(-\mathbf{\Phi}^{-1} \cdot \partial_{x_j}\bigr(\mathbf{\Phi}^{-1}\bigl) \cdot \mathbf{\Phi}^{-1}\Biggl)_{k,k}.
\end{equation}
and the derivative of $\mathbf{\Phi}$ as
\begin{align}
	\partial_{x_j}\bigr(\mathbf{\Phi}^{-1}\bigl) &= \partial_{x_j} \biggr(\mathbf{X}\mathbf{H}^H\frac{1}{\sigma^2}\mathbf{H} \mathbf{X} + \mathbf{I_M} \biggl)\\
	&= \partial_{x_j} \biggr(\mathbf{X}\mathbf{H}^H\frac{1}{\sigma^2}\mathbf{H} \mathbf{X}\biggl)\\
	&= \mathbf{E}_{j,j}\mathbf{H}^H\frac{1}{\sigma^2}\mathbf{HX} + \mathbf{XH}^H\frac{1}{\sigma^2}\mathbf{HE}_{j,j}
\end{align}
and finally
\begin{equation}
	\begin{aligned}
		\partial_{X_j}(R_k) = \frac{1}{2}\Biggr(\sum_k{(R_k-R_{tgt})^2}\Biggl)^{-1/2}\cdot\sum_k{\Biggr[2(R_k-R_{tgt})\cdot\frac{1}{\log{2}\cdot\mathbf{\Phi}_{k,k}^{-1}}\cdot\\\cdot\Biggr(-\mathbf{\Phi}^{-1} \cdot\Biggr( \mathbf{E}_{j,j}\mathbf{H}^H\frac{1}{\sigma^2}\mathbf{HX} + \mathbf{XH}^H\frac{1}{\sigma^2}\mathbf{HE}_{j,j}\Biggl) \cdot \mathbf{\Phi}^{-1}\Biggl)_{k,k}\Biggl]}.
	\end{aligned}
\end{equation}

\input{chapters/03-results}

%=== Appendix ===================================================
%\settocdepth{chapter} % Appendix nur Kapitelweise in Contents
\appendix
%\input{chapters/a-cdcontents}

%=== Bibliography ===============================================
%\nocite{*}
\cleardoublepage
\addcontentsline {toc}{chapter}{Bibliography} 
\bibliographystyle{IEEEtran}
\bibliography{IEEEabrv,report}
\cleardoublepage

%=== Index ======================================================
%\addcontentsline{toc}{chapter}{Index}
%\printindex 
%\cleardoublepage
\end{document}

%%%%%%%%%%%%%%%%%%%%%%%%%%%%%%%%%%%%%%%%%%%%%%%%%%%%%%%%%%%%%%%%%%%%
%%%  the  E N D  %%%%%%%%%%%%%%%%%%%%%%%%%%%%%%%%%%%%%%%%%%%%%%%%%%%
%%%%%%%%%%%%%%%%%%%%%%%%%%%%%%%%%%%%%%%%%%%%%%%%%%%%%%%%%%%%%%%%%%%%
